% Full answers PDF: MCQ answers + worked explanations per option, then official MS clips.
% Compiler: XeLaTeX.
\documentclass[12pt,a4paper]{article}
\usepackage[margin=18mm,headheight=14pt]{geometry}
\usepackage{fontspec}
\usepackage[version=4]{mhchem}
\usepackage{graphicx}
\usepackage{xcolor}
\usepackage{booktabs}
\usepackage{fancyhdr}
\usepackage{enumitem}
\usepackage{needspace}

\setmainfont{Helvetica Neue}

\pagestyle{fancy}
\fancyhf{}
\fancyhead[L]{\small 3.2 Ionic bonding --- answers}
\fancyhead[R]{\small CIE 9701}
\fancyfoot[C]{\thepage}
\renewcommand{\headrulewidth}{0.3pt}

\setlength{\parindent}{0pt}
\setlength{\parskip}{0.35em}

\newcommand{\msblock}[2]{%
  \par\vspace{0.8em}%
  \noindent\begin{minipage}{\textwidth}%
  {\normalsize\bfseries #1}\par\vspace{0.25em}%
  \noindent\includegraphics[width=\textwidth,keepaspectratio]{#2}\par
  \end{minipage}\par
}

\newcommand{\mcq}[4]{%
  \par\needspace{5\baselineskip}\vspace{0.7em}%
  {\normalsize\bfseries #1.\quad #2}%
  \hfill{\footnotesize\color{gray}#3}\par\vspace{0.15em}%
  {\small #4}\par
}

\newcommand{\why}[2]{\par\needspace{2.2\baselineskip}\vspace{0.1em}{\small\textbf{#1} #2}\par}

\begin{document}

\begin{center}
{\large\bfseries 3.2 Ionic bonding homework --- answers}\\[0.2em]
{\small CIE 9701 AS Chemistry \quad·\quad 7 multiple-choice + 10 written marks}
\end{center}

{\small\bfseries Do not issue this file to students.}
\hrule
\vspace{0.4em}

\section*{Section A --- Multiple choice}

{\normalsize\bfseries Quick answers}

\vspace{0.3em}
{\small
\begin{tabular}{@{}ll@{\hspace{2.6em}}ll@{\hspace{2.6em}}ll@{}}
\toprule
Q & Ans. & Q & Ans. & Q & Ans.\\
\midrule
A1 & \textbf{B} & A4 & \textbf{C} & A7 & \textbf{C}\\
A2 & \textbf{C} & A5 & \textbf{D} & & \\
A3 & \textbf{B} & A6 & \textbf{B} & & \\
\bottomrule
\end{tabular}}

\vspace{0.6em}
{\normalsize\bfseries Worked explanations}

\mcq{A1}{B}{2021 MJ Paper 13 Q6}{%
Ionic bonding requires the \emph{transfer} of electrons from a metal atom to a
non-metal atom --- two \emph{different} elements. An element contains only one
type of atom, so it cannot produce oppositely charged ions; ionic bonding is
therefore \emph{never} found in elements.}

\why{A:}{covalent --- \ce{H2}, \ce{Cl2}, diamond (giant covalent).}

\why{C:}{metallic --- Na, Fe, Cu.}

\why{D:}{van der Waals' forces \emph{are} found in elements --- between
\ce{I2} molecules or between noble-gas atoms.}

\mcq{A2}{C}{2025 ON Paper 11 Q16}{%
The radii identify the elements. A metal loses electrons, so its ionic radius is
\emph{smaller} than its atomic radius; a non-metal gains electrons, so its ionic
radius is \emph{larger}.
X: 0.118 $\rightarrow$ 0.053 nm --- a metal cation, matches Al (Al 0.118,
\ce{Al^3+} 0.054).
Y: 0.099 $\rightarrow$ 0.180 nm --- a non-metal anion, matches Cl (Cl 0.099,
\ce{Cl-} 0.181).
Z: 0.160 $\rightarrow$ 0.072 nm --- a metal cation, matches Mg (Mg 0.160,
\ce{Mg^2+} 0.072).
So X = Al, Y = Cl and Z = Mg.}

\why{A:}{X is aluminium, a metal that conducts electricity; Y is chlorine, a
non-metal that does not. X has the \emph{higher} conductivity, not lower.}

\why{B:}{Y is chlorine (a simple molecular gas, mp $-$101\,\textdegree C); X is
aluminium (mp 660\,\textdegree C). Y has the \emph{much lower} melting point.}

\why{C (correct):}{Y (Cl) and Z (Mg) are a non-metal and a metal, so they form
the ionic compound \ce{MgCl2}.}

\why{D:}{Z is magnesium, a metal --- it forms ions by \emph{losing} electrons,
not gaining them.}

\mcq{A3}{B}{2025 MJ Paper 11 Q8}{%
The key on the paper defines the symbols: \textbf{Na} and \textbf{Cl} stand for
the \emph{atoms}, while \textbf{+} and \textbf{--} stand for the \emph{ions}
(\ce{Na+} and \ce{Cl-}). A sodium chloride crystal is a giant ionic lattice of
\emph{ions} in a regular alternating arrangement, so the correct diagram is the
one built from + and -- symbols (option B).}

\why{A:}{labels the particles Na and Cl, i.e.\ \emph{atoms} (see the key). A
sodium chloride lattice contains \ce{Na+} and \ce{Cl-} ions, not atoms.}

\why{C / D:}{also use ion symbols, but the ions are not arranged in the correct
regular alternating (rock-salt) lattice of sodium chloride --- e.g.\ the packing
and relative positions of the two ions are wrong.}

\mcq{A4}{C}{2023 MJ Paper 11 Q6}{%
In the rock-salt structure each \ce{Na+} ion is surrounded by \textbf{6}
\ce{Cl-} ions (octahedral), and each \ce{Cl-} ion is surrounded by \textbf{6}
\ce{Na+} ions. The coordination numbers are therefore 6 and 6.}

\why{A / B:}{give 4 for one of the ions. A coordination number of 4 would be a
different structure (e.g.\ zinc blende), not NaCl.}

\why{D:}{gives 8 for the sodium ions. 8:8 is the caesium chloride structure;
\ce{Cs+} is much larger than \ce{Na+} so \ce{CsCl} packs differently.}

\mcq{A5}{D}{2025 ON Paper 13 Q18}{%
\ce{P4O10} is a simple molecular oxide: discrete \ce{P4O10} molecules held
together by weak van der Waals' forces, so it has a simple structure.
\ce{MgO}, \ce{Al2O3} and \ce{SiO2} all have \emph{giant} structures (giant ionic
for MgO and \ce{Al2O3}, giant covalent/macromolecular for \ce{SiO2}).}

\why{A:}{\ce{MgO} is a giant ionic lattice --- high melting point, conducts when
molten.}

\why{B:}{\ce{Al2O3} is a giant ionic lattice (with some covalent character).}

\why{C:}{\ce{SiO2} is a giant covalent (macromolecular) solid --- this is the
common trap, because it looks like a simple oxide of a non-metal but is
macromolecular.}

\mcq{A6}{B}{2025 MJ Paper 14 Q12}{%
Sodium ethanoate, \ce{CH3COONa}, contains \emph{ionic} bonding (between the
\ce{Na+} ion and the \ce{CH3COO-} ion) and \emph{covalent} bonding (the C--H,
C--C, C--O and \ce{C=O} bonds inside the ethanoate ion). Neither of these is a
coordinate (dative covalent) bond, so B is the answer.}

\why{A:}{\ce{Al2Cl6} contains covalent bonds \emph{and} coordinate bonds (the
Cl $\rightarrow$ Al dative bonds in the dimer), so it is excluded.}

\why{C:}{\ce{MgCl2} is ionic only --- it has no covalent bonding.}

\why{D:}{\ce{NH4Cl} contains ionic bonding plus covalent N--H bonds \emph{and} a
coordinate bond in the \ce{NH4+} ion (the bond formed when the lone pair on N
donates to \ce{H+}), so it is excluded.}

\mcq{A7}{C}{2021 ON Paper 12 Q35}{%
(1) \textbf{False} --- \ce{SiO2} is a giant covalent solid (mp about
1710\,\textdegree C) while \ce{SiCl4} is simple molecular (mp
$-$68\,\textdegree C); \ce{SiO2} has the much higher melting point.
(2) \textbf{True} --- \ce{SiCl4} is hydrolysed by water (\ce{SiCl4 + 2H2O ->
SiO2 + 4HCl}); \ce{NaCl} simply dissolves in water without reaction.
(3) \textbf{True} --- \ce{SiCl4} is simple molecular (weak van der Waals'
forces) so its melting point is far lower than that of \ce{NaCl}, a giant ionic
lattice.
Correct statements: 2 and 3 $\rightarrow$ C.}

\why{A / B:}{both include statement 1, which is false.}

\why{D:}{includes only statement 1 (false) and drops the two true statements.}

\newpage

\section*{Section B --- Mark scheme clips}

\msblock{B1(a)  [2]}{cie-9701-2022-on-p22-q1a-ms.png}

\msblock{B2(a)(i)  [1]}{cie-9701-2021-mj-p22-q2a-i-ms.png}

\msblock{B2(a)(ii)  [2]}{cie-9701-2021-mj-p22-q2a-ii-ms.png}

\msblock{B3(a)(i)  [1]}{cie-9701-2025-mj-p21-q1a-i-ms.png}

\msblock{B3(a)(iii)  [1]}{cie-9701-2025-mj-p21-q1a-iii-ms.png}

\msblock{B4(a)(i)  [1]}{cie-9701-2022-on-p23-q3a-i-ms.png}

\msblock{B5(c)(ii)  [2]}{cie-9701-2025-mj-p23-q2c-ii-ms.png}

\end{document}
